%!TEX root = ../template.tex
%%%%%%%%%%%%%%%%%%%%%%%%%%%%%%%%%%%%%%%%%%%%%%%%%%%%%%%%%%%%%%%%%%%%
%% abstract-pt.tex
%% NOVA thesis document file
%%
%% Abstract in Portuguese
%%%%%%%%%%%%%%%%%%%%%%%%%%%%%%%%%%%%%%%%%%%%%%%%%%%%%%%%%%%%%%%%%%%%

\typeout{NT FILE abstract-pt.tex}%

Os Conflict-Free Replicated Data Types (CRDTs) representam um conceito fundamental para sistemas distribuídos modernos,
permitindo atualizações concorrentes em múltiplas réplicas ao mesmo tempo que garantem Strong Eventual Consistency sem
coordenação centralizada. Apesar destas garantias, desenhar CRDTs corretamente continua a ser um desafio, e a verificação
formal desempenha um papel importante na prevenção de erros que comprometam a integridade dos dados. As ferramentas de
verificação automática baseadas em solvers SMT, como o VeriFx, permitem a rápida validação experimental de propriedades
algébricas de convergência, mas enfrentam limitações significativas ao raciocinar sobre definições recursivas,
propriedades indutivas e invariantes de estado globais. Plataformas dedutivas como o Why3 ultrapassam estas limitações
através de contratos e invariantes explícitos, ao custo de um esforço de especificação muito superior por parte do
programador.

Esta dissertação responde a este compromisso com uma framework integrada para o desenvolvimento modular e formal de CRDTs.
Ao invés de traduzir diretamente entre o VeriFx e o Why3, a framework introduz uma linguagem própria para especificar
CRDTs baseados em estado e em operações. Depois de validada através de análise léxica, sintática e de tipos sobre uma
árvore de sintaxe abstrata comum, uma única especificação é automaticamente compilada tanto para VeriFx, visando a prova
automática de propriedades de convergência através do solver Z3, como para WhyML, visando a verificação dedutiva da
correção semântica e dos invariantes globais. Este desenho oculta do programador a complexidade sintática e de modelação
de ambos os ecossistemas alvo, assegurando que os dois esforços de verificação permanecem consistentes entre si, já que
ambos derivam da mesma especificação. A framework foi avaliada em dez CRDTs, incluindo contadores, conjuntos, um mapa
genérico e o Replicated Growable Array, reduzindo o esforço de especificação manual exigido, ao mesmo tempo que gera
implementações que respeitam as convenções específicas de cada ferramenta. Em particular, o caso de estudo do Replicated
Growable Array expõe um cenário em que a prova dedutiva no Why3 é bem-sucedida enquanto a prova totalmente automática no
VeriFx falha, ilustrando a vantagem de combinar verificação automática e dedutiva em vez de depender de qualquer uma
isoladamente.

% E agora vamos fazer um teste com uma quebra de linha no hífen a ver se a \LaTeX\ duplica o hífen na linha seguinte se usarmos \verb+"-+… em vez de \verb+-+.
%
% zzzz zzz zzzz zzz zzzz zzz zzzz zzz zzzz zzz zzzz zzz zzzz zzz zzzz zzz zzzz comentar"-lhe zzz zzzz zzz zzzz
%
% Sim!  Funciona! :)

% Palavras-chave do resumo em Português
% \begin{keywords}
% Palavra-chave 1, Palavra-chave 2, Palavra-chave 3, Palavra-chave 4
% \end{keywords}
\keywords{
  CRDTs \and
  Verificação Formal \and
  Why3 \and
  VeriFx
}
% to add an extra black line
